% Generado por Experimentos/generar_resultados_latex.py desde complementario.
% No editar a mano: volver a ejecutar el script.

\begin{table}[ht]
\centering
\footnotesize
\begin{tabular}{c c c c}
\hline
Ciclo & $\alpha=0$ & $\alpha=0.5$ & $\alpha=1$ \\
\hline
0 ($M_0$) & \multicolumn{3}{c}{$0.8827$ [$0.8752$; $0.8901$]} \\
1 & $0.8863$ [$0.8799$; $0.8949$] & $0.8859$ [$0.8739$; $0.8942$] & $0.8864$ [$0.8777$; $0.8939$] \\
2 & $0.8898$ [$0.8861$; $0.8942$] & $0.8954$ [$0.8909$; $0.8982$] & $0.8918$ [$0.8825$; $0.8980$] \\
3 & $0.8886$ [$0.8816$; $0.9009$] & $0.8873$ [$0.8815$; $0.8904$] & $0.8994$ [$0.8934$; $0.9063$] \\
4 & $0.8945$ [$0.8883$; $0.9059$] & $0.8948$ [$0.8941$; $0.8957$] & $0.9017$ [$0.8923$; $0.9066$] \\
5 & $0.8903$ [$0.8862$; $0.8974$] & $0.8984$ [$0.8925$; $0.9018$] & $0.9029$ [$0.9017$; $0.9039$] \\
\hline
\end{tabular}
\caption{Dispersión semántica: media y rango [mínimo; máximo] de las tres cadenas por ciclo y proporción sintética $\alpha$. $M_0$ es compartido por las tres proporciones dentro de cada repetición. Salvo la perplejidad, calculada con los 128 contextos de prueba, las métricas corresponden al análisis complementario \textit{post hoc} con los 90 contextos en los que todas las etapas alcanzan $W=30$ palabras.}
\label{tab:medias_semantic_dispersion}
\end{table}
